\documentclass[../../main.tex]{subfiles}% 注意这里的文件路径不能用 ./main.tex ,否则用latexmk编译子文件会报错
\graphicspath{{\subfix{./image/}}{\subfix{../../image/}}} % 指定图片引用路径,后续可以直接使用图片文件名
% 如果图片存放在上级目录,则使用 ../../image/ 作为备用路径

% 例如:
% \begin{figure}[H]
% \centering
% \includegraphics[width=0.8\textwidth]{图.png}
% \caption{}
% \label{figure:图}
% \end{figure}
% 注意:上述\label{}一定要放在\caption{}之后,否则引用图片序号会只会显示??.

\begin{document}

\section{对椭圆方程的运用}

\begin{examples}
设$\Omega\subseteq\mathbb{R}^n$为有界、具有光滑边界的开区域，$U(x)\in C(\overline{\Omega})$，$f\in L^2(\Omega)$，研究椭圆型Dirichlet边值问题
\begin{align}
\begin{cases}
-\Delta u+U(x)u=f(x) &(x\in\Omega),\\
u|_{\partial\Omega}=0.
\end{cases}
\label{eq:elliptic_dirichlet}
\end{align}
\end{examples}

\begin{definition}
称$u\in H_0^1(\Omega)$为边值问题\eqref{eq:elliptic_dirichlet}的\textbf{弱解}，若对任意$v\in H_0^1(\Omega)$，成立
\begin{align}
\int_{\Omega}\left(\nabla u\cdot\nabla v+U(x)uv\right)\mathrm{d}x=\int_{\Omega}fv\mathrm{d}x.
\label{eq:weak_sol}
\end{align}
\end{definition}

\begin{remark}
若$u\in H^2(\Omega)$满足\eqref{eq:elliptic_dirichlet}，则$u$必为其弱解；反之，\eqref{eq:elliptic_dirichlet}的弱解必属于$H^2(\Omega)$，故仅需研究弱解的存在性。
\end{remark}

\begin{definition}
记$\lambda_0\triangleq\max\limits_{x\in\overline{\Omega}}|U(x)|$，在$H_0^1(\Omega)$上定义等价范数
\begin{align*}
\square u\square=\left[\int_{\Omega}|\nabla u|^2\mathrm{d}x+\int_{\Omega}\left(U(x)+\lambda_0\right)u^2(x)\mathrm{d}x\right]^{\frac12},
\end{align*}
同时定义内积
\begin{align}
(u,v)_{\lambda_0}=\int_{\Omega}\nabla u\cdot\nabla v\mathrm{d}x+\int_{\Omega}\left(U(x)+\lambda_0\right)u(x)v(x)\mathrm{d}x,\quad \forall u,v\in H_0^1(\Omega).
\label{eq:inner_lambda0}
\end{align}
\end{definition}

\begin{proposition}
\begin{enumerate}[(1)]
\item 存在常数$m,M>0$，使得
\begin{align}
m\|u\|_1\leqslant\square u\square\leqslant M\|u\|_1,\quad \forall u\in H_0^1(\Omega),
\label{eq:equiv_norm}
\end{align}
其中$\|\cdot\|_1$为$H_0^1(\Omega)$的标准范数，故$\square\cdot\square$是$H_0^1(\Omega)$的等价范数，$H_0^1(\Omega)$在该范数下为$B$空间；
\item $H_0^1(\Omega)$在内积\eqref{eq:inner_lambda0}下为Hilbert空间。
\end{enumerate}
\end{proposition}
\begin{proof}
由Poincaré不等式可直接推得\eqref{eq:equiv_norm}，结合等价范数性质与内积诱导范数的完备性，可得结论。

\end{proof}

\begin{definition}
\begin{enumerate}[(1)]
\item 对任意$u\in L^2(\Omega)$，由Riesz表示定理，存在唯一$w\in H_0^1(\Omega)$使得
\begin{align}
\int_{\Omega}uv\mathrm{d}x=(w,v)_{\lambda_0},\quad \forall v\in H_0^1(\Omega),
\label{eq:riesz_rep}
\end{align}
定义算子$K_{\lambda_0}:L^2(\Omega)\to H_0^1(\Omega)$满足$w=K_{\lambda_0}u$；
\item 记$\iota:H_0^1(\Omega)\to L^2(\Omega)$为自然嵌入算子。
\end{enumerate}
\end{definition}

\begin{proposition}
\begin{enumerate}[(1)]
\item $K_{\lambda_0}$为连续线性算子，且存在常数$c>0$，满足$\square K_{\lambda_0}u\square\leqslant c\|u\|,\ \forall u\in L^2(\Omega)$；
\item 由Rellich定理，$\iota$为紧算子，故复合算子$K_{\lambda_0}\iota:H_0^1(\Omega)\to H_0^1(\Omega)$为紧对称算子。
\end{enumerate}
\end{proposition}
\begin{proof}
(1) 由Poincaré不等式，$\left|\int_{\Omega}uv\mathrm{d}x\right|\leqslant\|u\|\|v\|\leqslant c\|u\|\square v\square$，结合Riesz表示定理可得算子的连续性与有界性；
(2) 紧算子与连续算子的复合仍为紧算子，内积的对称性可推得算子对称性。

\end{proof}

\begin{proposition}
弱解方程\eqref{eq:weak_sol}等价于$H_0^1(\Omega)$中的算子方程
\begin{align}
\left(I-\lambda_0 K_{\lambda_0}\iota\right)u=K_{\lambda_0}f.
\label{eq:operator_eq}
\end{align}
\end{proposition}
\begin{proof}
将\eqref{eq:weak_sol}变形为
\begin{align*}
\int_{\Omega}\left(\nabla u\cdot\nabla v+\left(U+\lambda_0\right)uv\right)\mathrm{d}x=\int_{\Omega}fv\mathrm{d}x+\lambda_0\int_{\Omega}uv\mathrm{d}x,
\end{align*}
即$(u,v)_{\lambda_0}=(K_{\lambda_0}f,v)_{\lambda_0}+\lambda_0(K_{\lambda_0}\iota u,v)_{\lambda_0},\ \forall v\in H_0^1(\Omega)$，
由内积的非退化性即得\eqref{eq:operator_eq}。

\end{proof}

\begin{theorem}
设$\Omega\subseteq\mathbb{R}^n$为有界光滑区域，$U\in C(\overline{\Omega})$，$f\in L^2(\Omega)$，对边值问题\eqref{eq:elliptic_dirichlet}：
\begin{enumerate}[(1)]
\item 若其对应的齐次问题仅有零弱解，则对任意$f\in L^2(\Omega)$，原问题存在唯一弱解；
\item 若齐次问题存在非零弱解，则齐次问题至多存在有限个线性无关弱解$\{\varphi_1,\varphi_2,\cdots,\varphi_n\}$，此时原问题有解当且仅当
\begin{align*}
\int_{\Omega}f\varphi_i\mathrm{d}x=0,\quad i=1,2,\cdots,n,
\end{align*}
且解空间维数为$n$。
\end{enumerate}
\end{theorem}
\begin{proof}
由Riesz-Fredholm定理，算子方程\eqref{eq:operator_eq}有解当且仅当$K_{\lambda_0}f\perp N\left(I-\lambda_0 K_{\lambda_0}\iota\right)$。
验证零空间等价性：
\begin{align*}
u\in N\left(I-\lambda_0 K_{\lambda_0}\iota\right)&\iff u=\lambda_0 K_{\lambda_0}\iota u\\
&\iff (u,v)_{\lambda_0}=\lambda_0\int_{\Omega}uv\mathrm{d}x,\ \forall v\in H_0^1(\Omega)\\
&\iff \int_{\Omega}\left(\nabla u\cdot\nabla v+U(x)uv\right)\mathrm{d}x=0,\ \forall v\in H_0^1(\Omega),
\end{align*}
即$u$为齐次问题的弱解。
设齐次问题线性无关弱解为$\{\varphi_1,\cdots,\varphi_n\}$，则
\begin{align*}
(K_{\lambda_0}f,\varphi_i)_{\lambda_0}=0\iff\int_{\Omega}f\varphi_i\mathrm{d}x=0,\quad i=1,\cdots,n,
\end{align*}
结合Fredholm择一性，有限维零空间与正交条件可得定理结论。

\end{proof}

\begin{proposition}
考虑对应特征值边值问题
\begin{align}
\begin{cases}
-\Delta u+U(x)u=\lambda u &(x\in\Omega),\\
u|_{\partial\Omega}=0.
\end{cases}
\label{eq:eigenvalue_pb}
\end{align}
特征值问题\eqref{eq:eigenvalue_pb}等价于算子方程
\begin{align*}
\left(I-\left(\lambda+\lambda_0\right)K_{\lambda_0}\iota\right)u=0.
\end{align*}
\end{proposition}
\begin{proof}
\eqref{eq:eigenvalue_pb}的弱形式为
\begin{align*}
\int_{\Omega}\left(\nabla u\cdot\nabla v+Uuv\right)\mathrm{d}x=\lambda\int_{\Omega}uv\mathrm{d}x,\quad\forall v\in H_0^1(\Omega),
\end{align*}
整理得
\begin{align*}
(u,v)_{\lambda_0}=\left(\lambda+\lambda_0\right)\int_{\Omega}uv\mathrm{d}x=\left(\lambda+\lambda_0\right)(K_{\lambda_0}\iota u,v)_{\lambda_0},
\end{align*}
即$\left(I-\left(\lambda+\lambda_0\right)K_{\lambda_0}\iota\right)u=0$。

\end{proof}

\begin{theorem}
边值问题\eqref{eq:eigenvalue_pb}的特征值均为实数，可列个且满足$\lambda_1\leqslant\lambda_2\leqslant\cdots\leqslant\lambda_j\leqslant\cdots$，$\lambda_j\to\infty\ (j\to\infty)$；对应特征函数构成$H_0^1(\Omega)$的完备正交集。
\end{theorem}
\begin{proof}
由Riesz-Schauder理论与Hilbert-Schmidt定理，紧对称算子$K_{\lambda_0}\iota$的非零谱$\sigma(K_{\lambda_0}\iota)\setminus\{0\}$为实数，可列个，记为$\{\mu_j\}$，且$\mu_j\to0\ (j\to\infty)$。
由$(K_{\lambda_0}\iota u,u)_{\lambda_0}>0,\ \forall u\neq0$，得$\mu_j>0$，不妨设$\mu_1\geqslant\mu_2\geqslant\cdots\geqslant\mu_j\geqslant\cdots>0$。
由算子方程关系，特征值满足$\lambda_j=\frac{1}{\mu_j}-\lambda_0$，故$\lambda_1\leqslant\lambda_2\leqslant\cdots$且$\lambda_j\to\infty$。
紧对称算子的特征函数系完备正交，故原问题特征函数构成$H_0^1(\Omega)$的完备正交集。

\end{proof}

\begin{proposition}
特征值$\lambda_j$具有极小极大刻画：
\begin{align*}
\lambda_j=\sup_{E_{j-1}}\inf_{\substack{u\in E_{j-1}^\perp\\u\neq\theta}}\frac{\int_{\Omega}\left(|\nabla u|^2+Uu^2\right)\mathrm{d}x}{\int_{\Omega}u^2\mathrm{d}x},
\end{align*}
其中$E_{j-1}$为$H_0^1(\Omega)$中任意$j-1$维闭线性子空间。
\end{proposition}
\begin{proof}
紧对称算子特征值满足极小极大原理：
\begin{align*}
\mu_j=\inf_{E_{j-1}}\sup_{\substack{u\in E_{j-1}^\perp\\u\neq\theta}}\frac{(K_{\lambda_0}\iota u,u)_{\lambda_0}}{(u,u)_{\lambda_0}},
\end{align*}
结合$\lambda_j=\frac{1}{\mu_j}-\lambda_0$，代入内积与范数定义即得结论。

\end{proof}

\begin{example}
设 $a_i(x) \in C^1(\Omega)(i=1,2,\cdots,n)$，$U(x) \in C(\overline{\Omega})$，其中$\Omega$是$\mathbb{R}^n$中的边界光滑的有界开区域，讨论下列边值问题：
\begin{align*}
\begin{cases}
-\Delta u + \sum\limits_{i=1}^{n} \partial_{x_i}(a_i(x)u) + U(x)u = f(x) & (x \in \Omega), \\
u\big|_{\partial\Omega} = 0.
\end{cases}
\end{align*}
提示：应用Lax‑Milgram定理（定理2.3.18）。
\end{example}

\begin{example}
在上题中，讨论下列特征值问题：
\begin{align*}
\begin{cases}
-\Delta u + \sum\limits_{i=1}^{n} \partial_{x_i}(a_i(x)u) + U(x)u = \lambda u & (x \in \Omega), \\
u\big|_{\partial\Omega} = 0.
\end{cases}
\end{align*}
\end{example}


\end{document}